\documentclass[letterpaper,11pt]{article}
\usepackage[top=0.8in, bottom=0.8in, left=0.9in, right=0.9in]{geometry}
\usepackage[hidelinks]{hyperref}
\usepackage{lmodern}
\usepackage{parskip}
\usepackage{microtype}
\setlength{\parskip}{6pt}
\setlength{\parindent}{0pt}
\begin{document}
\begin{flushleft}
\textbf{\large Ajay Venkatesh} \\
Brooklyn, NY \\
+1 516 585 9013 \\
\href{mailto:av3855@nyu.edu}{av3855@nyu.edu}
\end{flushleft}
\vspace{10pt}
May 21, 2026
\vspace{6pt}

Hiring Manager \\
ASSA ABLOY \\
Phoenix, AZ

\vspace{10pt}

Dear Hiring Manager,

My name is Ajay Venkatesh. I'm finishing my M.S. in Computer Engineering at NYU, with production software experience at PwC in Bengaluru, a summer SDE internship at Amazon, and ongoing on-campus work at NYU's Novel AI Technologies. I'm writing about the Junior Software Engineer role with the Openings Studio team in Phoenix.

The role's mix of \textbf{Java}, \textbf{.NET}, and \textbf{native C++} sits at the intersection of where I've been compounding. At Amazon I engineered backend services in \textbf{Java} (Coral, Smithy, Dagger, CBOR) with low-level design patterns and unit and integration testing via \textbf{Mockito}, alongside a \textbf{React 18} frontend with optimized API integration. The work was grounded in \textbf{object-oriented design}, careful API contracts, and code-review discipline across a small team.

On the \textbf{.NET} and Microsoft stack: at PwC I owned a distributed backend on \textbf{Microsoft Azure} integrating with \textbf{Microsoft Dynamics CRM} over \textbf{SQL Server}, with event-driven processing, fault-tolerant retries, and parallelized batch ingestion. I worked daily inside \textbf{Visual Studio} and \textbf{Azure DevOps} on a multi-language codebase, sized enhancement requests with architects and PMs, and shipped under a SDLC governance model with strict review gates. \textbf{C++} is where I first learned object-oriented programming in graduate coursework, and I've stayed comfortable with it through \textbf{Data Structures} and \textbf{System Design} at NYU.

At NYU I shipped a customer-facing analytics platform on \textbf{React, Node.js, REST APIs} deployed via \textbf{Docker} on \textbf{EC2}, partnered with non-technical stakeholders to translate requirements, and learned how to write code that has to live in someone else's tools downstream. That integration mindset, building things that interface with third-party design tooling, plug-ins, and a multi-vendor surface, is what Openings Studio sounds like.

Stack-wise: \textbf{Java}, \textbf{.NET}, \textbf{C++}, \textbf{Python}, \textbf{JavaScript / TypeScript}, \textbf{SQL Server / PostgreSQL / MySQL}, with \textbf{Visual Studio}, \textbf{Eclipse}, \textbf{IntelliJ}, \textbf{Git}, and \textbf{Azure DevOps} as daily tools.

What pulls me toward ASSA ABLOY is the product itself. \textbf{Openings Studio} sits in front of real construction projects, real hardware specifications, and real installations in airports, hospitals, and hotels. Software that compounds into physical buildings is where engineering judgment has consequences you can walk through, and that kind of grounded work is what I'd like to build my career around.

I'd welcome the chance to discuss further. Thanks for considering my application.

\vspace{12pt}
Sincerely, \\
Ajay Venkatesh

\end{document}
